%%%%%%%%%%%%%%%%%%%%%%%%%%%%%%%%%%%%%%%%%%%%%%%%%%%%%%%%%%%%%%%%%%%%%%%%%%%%%%%%%%%%
%%%%%%%%%%%%%%%%%%%%%%%%%%%%%%%%%%%%%%%%%%%%%%%%%%%%%%%%%%%%%%%%%%%%%%%%%%%%%%%%%%%%
%
% Some functional relations derived from the Lindel\"of-Wirtinger expansion
% of the Lerch transcendent function,
% by
% L. M. Navas, F. J. Ruiz and J. L. Varona
%
% Date of first version: January 11, 2011
% Date of last modification: July 28, 2013
%
%%%%%%%%%%%%%%%%%%%%%%%%%%%%%%%%%%%%%%%%%%%%%%%%%%%%%%%%%%%%%%%%%%%%%%%%%%%%%%%%%%%%
%%%%%%%%%%%%%%%%%%%%%%%%%%%%%%%%%%%%%%%%%%%%%%%%%%%%%%%%%%%%%%%%%%%%%%%%%%%%%%%%%%%%

%\documentclass[reqno]{amsart}
\documentclass[11pt,a4paper,reqno]{amsart}

%%%%%%%%%%%%%%%%%%%%%%%%%%%%%%%%%%%%%%%%%%%%%%%%%%%%%%%%%%%%%%%%%%%%%%%%%%%%%%%%%%%


\usepackage{url}
\usepackage[colorlinks,linkcolor=blue,citecolor=blue]{hyperref}

% \usepackage[hyperpageref]{backref} % For drafts

%\date{July 28, 2013}
\renewcommand{\datename}{\textbf{THIS PAPER HAS PUBLISHED IN:}}
\date{\textit{Mathematics of Computation} \textbf{84} (2015), no.~292, 803--813.}


%%%%%%%%%%%%%%%%%%%%%%%%%%%%%%%%%%%%%%%%%%%%%%%%%%%%%%%%%%%%%%%%%%%%%%%%%%%%%%%%%%%%


\newcommand{\B}{\mathcal{B}}
\newcommand{\C}{\mathbb{C}}
\newcommand{\E}{\mathcal{E}}
\newcommand{\N}{\mathbb{N}}
\newcommand{\Q}{\mathbb{Q}}
\newcommand{\R}{\mathbb{R}}
\newcommand{\Z}{\mathbb{Z}}

%\renewcommand{\epsilon}{\varepsilon}

\renewcommand{\Re}{\operatorname{Re}}
\renewcommand{\Im}{\operatorname{Im}}


\theoremstyle{plain}
\newtheorem{theorem}{Theorem}%[section]
\newtheorem{prop}[theorem]{Proposition}
\newtheorem{cor}[theorem]{Corollary}
\newtheorem{lemma}{Lemma}

\theoremstyle{remark}
\newtheorem{remark}{Remark}

\numberwithin{equation}{section}
\allowdisplaybreaks

%%%%%%%%%%%%%%%%%%%%%%%%%%%%%%%%%%%%%%%%%%%%%%%%%%%%%%%%%%%%%%%%%%%%%%%%%%%%%%%%%%%%

\title[Functional relations from the Lindel\"of-Wirtinger expansion]%
{Some functional relations derived from the Lindel\"of-Wirtinger expansion of the Lerch transcendent function}


%\author{L.~M.~Navas, F.~J.~Ruiz and J.~L.~Varona}

\author[L.~M.~Navas]{Luis M. Navas}
\address{%
        Departamento de Matem\'aticas,
        Universidad de Salamanca,
        Plaza de la Merced 1-4,
        37008 Salamanca, Spain}
\email{navas@usal.es}

\author[F.~J.~Ruiz]{Francisco J. Ruiz}
\address{%
        Departamento de Matem\'aticas,
        Universidad de Zaragoza,
        Campus de la Plaza de San Francisco,
        50009 Zaragoza, Spain}
\email{fjruiz@unizar.es}

\author[J.~L.~Varona]{Juan L. Varona}
\address{%
        Departamento de Matem\'aticas y Computaci\'on,
        Universidad de La Rioja,
        Calle Luis de Ulloa s/n,
        26004 Logro\~no, Spain}
\email{jvarona@unirioja.es}
\urladdr{http://www.unirioja.es/cu/jvarona/welcome.html}

\thanks{The authors were supported by grant MTM2012-36732-C03-02 of the DGI}

\keywords{Lerch transcendent function, M\"obius inversion, Fourier series}

\subjclass[2010]{Primary 41A60; Secondary 11M35, 42A10}

%%%%%%%%%%%%%%%%%%%%%%%%%%%%%%%%%%%%%%%%%%%%%%%%%%%%%%%%%%%%%%%%%%%%%%%%%%%%%%%%%%%%




%%%%%%%%%%%%%%%%%%%%%%%%%%%%%%%%%%%%%%%%%%%%%%%%%%%%%%%%%%%%%%%%%%%%%%%%%%%%%%%%%%%%
\begin{document}
%%%%%%%%%%%%%%%%%%%%%%%%%%%%%%%%%%%%%%%%%%%%%%%%%%%%%%%%%%%%%%%%%%%%%%%%%%%%%%%%%%%%


%--------------------------------------
\begin{abstract}
The Lindel\"of-Wirtinger expansion of the Lerch transcendent function implies, as a limiting case, Hurwitz's formula for the eponymous zeta function. A generalized form of M\"obius inversion applies to the Lindel\"of-Wirtinger expansion and also implies an inversion formula for the Hurwitz zeta function as a limiting case. The inverted formulas involve the dynamical system of rotations of the circle and yield an arithmetical functional equation.
\end{abstract}
%--------------------------------------

%--------------------------------------
\maketitle
%--------------------------------------

%------------------
\section{Introduction}
\label{sec:intro}
%------------------


The Lerch transcendent function is given by the series
%
\begin{equation}
\label{eq:Lerch}
  \Phi(\lambda,s,z) = \sum_{k=0}^\infty\frac{\lambda^k}{(k+z)^s};
\end{equation}
%
see \cite[\S1.11, p.~27]{Bate} or \cite[\S25.14]{Ap-NIST}, for example.
Logarithms and complex powers are always assumed to be principal.
If $|\lambda| < 1$, then for any $s \in \C$, the series converges uniformly in $z$ over  $\C \setminus (-\infty,0]$, thus defining a holomorphic function of $z$ in this region. If $|\lambda| = 1$, then the series converges in this same region provided $\Re s > 1$. The value $\lambda = 0$ is considered trivial since it yields $\Phi(0,s,z) = z^{-s}$, and thus is usually excluded. There are multiple ways of defining analytic continuations of $\Phi$ in each parameter.

This function, defined by Mathias Lerch in 1887 in his paper \cite{Lerch},
includes as special cases of the parameters the Hurwitz and Riemann zeta functions
and the polylogarithms, among others, and has applications ranging from number theory to physics. It is often used to obtain functional identities; see, for instance, \cite{BaSa, Boy, GuiS, SC-book}.


One often extends the domain to $z \in \C\setminus \{0,-1,-2,\dots\}$ by including the branch discontinuity of the principal argument.
In addition, in this paper we shall only be considering $s \in \C$ with $\Re s < 0$, in which case the summand $(k+z)^{-s}$ (with $k \in \N$) in~\eqref{eq:Lerch} continuously extends to $z = -k$ by defining $0^{-s} = 0$. Thus for $\Re s < 0$ we can define $\Phi(\lambda,s,z)$ for all $z \in \C$.
For the parameter $s$, we shall often denote $\sigma = \Re s$.
Also, we shall use the notation $\C^{*} = \C \setminus \{0\}$. 

In light of the remarks made above, one concludes that given $|\lambda| < 1$ and $\Re s < 0$, the function $\Phi(\lambda,s,x)$ extends continuously as a function of $x \in [0,1]$, with $\Phi(\lambda,s,0) = \Phi(\lambda,s,1)$.
In particular, it may be expanded in a Fourier series on this interval. The resulting formula is classical and is the point of departure for what follows.


\begin{theorem}[Lindel\"of-Wirtinger expansion]
\label{th:L-W}
Let $\lambda$ and $s$ be complex parameters with $0<|\lambda|<1$ and $\Re s<0$. Then
%
\begin{equation}
\label{eq:L-W}
  \Phi(\lambda,s,x)
  = \lambda^{-x}\,\Gamma(1-s) \sum_{n\in\Z} (2\pi i n - \log\lambda)^{s-1}e^{2\pi i n x},
  \quad x\in[0,1].
\end{equation}
%
The convergence is uniform in~$x$.
\end{theorem}

\begin{proof}
Wirtinger's paper \cite{Wir} is one of the original sources. The result has been reproved often by various means. See for example~\cite{Bate}, for the traditional approach using complex analytic techniques, or~\cite{NRV-JAT2}, for a short proof using basic Fourier Analysis.
\end{proof}

\begin{remark}
Note that changing the branch of the logarithm in~\eqref{eq:L-W} shifts the summation index $n$, and hence does not affect the validity of the expansion, as long as the same branch is used for both $\log\lambda$ and~$\lambda^{-x}$.
\end{remark}


The special values $s = 1 - k$ (with $k \in \N$) give the \emph{Apostol-Bernoulli polynomials} $\B_{k}(x;\lambda)$ first defined in~\cite{Ap}:
%
\begin{equation}
\label{eq:A-B}
  \Phi(\lambda,1-k,x) = -\frac{\B_k(x;\lambda)}{k}.
\end{equation}
%
$\B_{k}(x;\lambda)$ is a polynomial in $x$ of degree $k-1$ over whose coefficients are rational functions in $\lambda$ having a unique pole at $\lambda = 1$.
By~\eqref{eq:L-W} we obtain their Fourier series
%
\begin{equation}
\label{eq:A-BFourier}
  \B_k(x;\lambda) = -\lambda^{-x}\,k!\, \sum_{n\in\Z} \frac{e^{2n\pi i x}}{(2n\pi i-\log\lambda)^k},
  \quad x\in[0,1],
\end{equation}
%
which is initially valid for $|\lambda|<1$, but is extended by analytic continuation to all $\lambda\neq 0,1$. In~\cite{NRV-MC}, \eqref{eq:A-BFourier} is proved directly using the algebraic properties of this polynomial family.

The structure of the paper is as follows. In Section~\ref{sec:LW-inv}, we look at the Lindel\"of-Wirtinger expansion in a manner analogous to Hurwitz's formula for the Hurwitz zeta function. The expansion~\eqref{eq:L-W} is separated into two parts involving a three-parameter function which specializes to the periodic zeta function. Next, we observe that one may apply to this new function an inversion formula for a certain generalized form of convolution, which contains the M\"obius inversion formula of analytic number theory as a special case (this general framework, discussed previously in~\cite{BNV}, is described in Section~\ref{sec:MobiusInvGen}). This yields an inverted form of the Lindel\"of-Wirtinger expansion. In Section~\ref{sec:HW}, by studying the logarithmic singularity of $\Phi$ at $\lambda = 1$, we show how to deduce Hurwitz's formula from the Lindel\"of-Wirtinger expansion and in Section~\ref{sec:HW-Inv}, using the same technique, we obtain the corresponding inverted form of Hurwitz's formula from the
inverted form of the Lindel\"of-Wirtinger expansion.
As special cases we obtain relations for the Apostol-Bernoulli polynomials.
Finally, in Section~\ref{sec:functionaleq}, for rational $x$, it is shown that the inverted form of Hurwitz's formula generalizes the functional equation of the Riemann zeta function. 


%------------------
\section{Generalized M\"obius inversion}
\label{sec:MobiusInvGen}
%------------------


In \cite{BNV} and \cite{BNV2}, an abstract framework for M\"obius inversion is established along with numerous examples. We will show that the Lindel\"of-Wirtinger formula (Theorem~\ref{th:L-W}) provides another application of this mechanism. In order to formulate the result, first we need to summarize its main features.
Although the results in \cite{BNV} are valid in a rather general setting, here we will only need the following special case.

Consider a dynamical system consisting of the natural numbers $\N$ acting on
a space $X$, in other words, a function $\varphi : \N \times X \to X$ such that $\varphi(1,x) = x$ and $\varphi(m,\varphi(n,x)) = \varphi(mn,x)$.
The action of $\N$ induces an action of the ring of complex-valued arithmetical functions on suitable spaces of functions $f : X \to \C$, where an arithmetical function $\alpha : \N \to \C$ acts on a function $f$ via
%
\begin{equation}
\label{eq:defAction}
	(\alpha \odot f)(x) = \sum_{n \in \N} \alpha(n) f(\varphi(n,x)),
	\quad
	x \in X.
\end{equation}
%
In other words, one has
%
\begin{equation*}
%\label{eq:Action}
	\alpha \odot (\beta \odot f) = (\alpha \ast \beta) \odot f
\end{equation*}
%
where $\ast$ denotes Dirichlet convolution of arithmetical functions, given by
\[
	(\alpha * \beta)(n) = \sum_{kl = n}	\alpha(k)\beta(l).
\]
Here, ``suitable'' means that the series~\eqref{eq:defAction} should converge fast enough. See~\cite{BNV} for the technical details. 

If $\alpha$ is invertible under convolution then, again under suitable convergence hypotheses, we have
%
\begin{equation*}
%\label{eq:inv}
	g = \alpha \odot f \implies f = \alpha^{-1} \odot g
\end{equation*}
%
where $\alpha^{-1}$ denotes the convolution inverse of $\alpha$. In particular, if $\alpha$ is completely multiplicative, its Dirichlet inverse is $\alpha^{-1} = \mu \alpha$, where $\mu$ is the M\"obius function. Thus we may expect in this case that
%
\begin{equation}
\label{eq:invCM}
	g = \alpha \odot f \implies f = (\mu \alpha) \odot g
\end{equation}
%
if the appropriate convergence conditions are satisfied.


%------------------
\section{M\"obius inversion of the Lindel\"of-Wirtinger formula}
\label{sec:LW-inv}
%------------------

We will apply the method of generalized M\"obius inversion outlined in Section~\ref{sec:MobiusInvGen} to the Lindel\"of-Wirtinger formula. This cannot be done directly with~\eqref{eq:L-W} as it stands, but rather after some manipulation of the series. Actually, we can separate the sum over all integers in~\eqref{eq:L-W} into a sum over positive and negative integers and invert each separately. For this purpose, consider the function defined by
%
\begin{equation}
\label{eq:L-defi}
	L(\lambda,s,x) = \sum_{n=1}^{\infty} (2\pi i n - \log\lambda)^{-s} e^{2\pi i nx}.
\end{equation}
%
This series is normally convergent for $x \in \R$ and on compact subsets of the domains defined by $\lambda \neq 0$, $\Re s > 1$. The principal branch of the logarithm is assumed throughout. Since $L$ is $1$-periodic in $x$ we may restrict to $x \in [0,1]$.


\begin{prop}
\label{prop:LW-sep}
With $L(\lambda,s,x)$ defined by~\eqref{eq:L-defi} and $0 < |\lambda| < 1$, $\Re s > 1$ and $x \in [0,1]$, we have
%
\begin{equation}
\label{eq:LW-sep}
	  \frac{\lambda^{x}}{\Gamma(s)}\Phi(\lambda,1-s,x) - (-\log \lambda)^{-s}
	= L(\lambda,s,x) + e^{\pi i s} L(\lambda^{-1},s,1-x).
\end{equation}
%
\end{prop}

\begin{proof}
Starting from the Lindel\"of-Wirtinger expansion~\eqref{eq:L-W}, with $s$ changed to $1-s$, and separating the term in the sum corresponding to $n=0$, gives
%
\begin{align*}
	   \frac{\lambda^{x}}{\Gamma(s)}\Phi(\lambda,1-s,x) - (-\log \lambda)^{-s}
	&= \sum_{n \in \Z \setminus \{0\}} (2 \pi i n - \log\lambda)^{-s} e^{2\pi i n x}.
\end{align*}
%
The sum over positive $n$ corresponds to $L(\lambda,s,x)$, while the sum over negative integers is
%
\begin{multline*}
	   \sum_{n=-\infty}^{-1} (2 \pi i n - \log\lambda)^{-s} e^{2\pi i n x}
	 = \sum_{n=1}^{\infty} (-2 \pi i n - \log\lambda)^{-s} e^{-2\pi i n x}
	\\
	 = e^{\pi i s} \sum_{n=1}^{\infty} (2 \pi i n + \log\lambda)^{-s} e^{2\pi i n (1-x)} 
	 = e^{\pi i s} L(\lambda^{-1},s,1-x),
\end{multline*}
%
since $\log \lambda^{-1} = -\log\lambda$ for $\lambda \notin (-\infty,0]$,
so that we obtain~\eqref{eq:LW-sep}.
\end{proof}


Next, we apply M\"obius inversion to~\eqref{eq:L-defi}.

\begin{prop}
\label{prop:L-inv}
For $\lambda \neq 0$, $\Re s > 1$ and $x \in [0,1]$ we have
%
\begin{equation}
\label{eq:L-inv}
	  (2 \pi i - \log \lambda)^{-s} e^{2\pi i x}
	= \sum_{n=1}^{\infty} \frac{\mu(n)}{n^{s}} L(\lambda^{1/n},s,\{nx\})
\end{equation}
% 
where $L(\lambda,s,x)$ is defined in~\eqref{eq:L-defi}, $\{x\}$ denotes the fractional part of a real number, and $\mu(n)$ is the M\"obius function, given by
\[
	  \mu(n)
	= \begin{cases}
	  0        & \text{if $n$ is not squarefree,}
	  \\
	  (-1)^{k} & \text{if $n$ is the product of $k$ distinct primes.}
	  \end{cases}
\]
\end{prop}

\begin{proof}
Consider the action of $\N$ on the space $\C^{*} \times [0,1)$ given by 
%
\begin{equation}
\label{eq:action}
	\varphi(n,(\lambda,x)) = \bigl(\lambda^{1/n}, \{nx\}\bigr). 
\end{equation}
%
It is straightforward to verify that $(\lambda^{a})^{b} = \lambda^{ab}$ holds for the principal power if $a,b \in (0,1]$, so that $(n,\lambda) \mapsto \lambda^{1/n}$ is an action on $\C^{*}$, and that $(n,x) \mapsto \{nx\}$ is an action on $[0,1]$. The action~\eqref{eq:action} is the direct product of these two.

Consider $\alpha(n) = n^{-s} : \N \to \C$, which is a completely multiplicative arithmetical function, and $f_{s} : \C^{*} \times [0,1] \to \C$ given by $f_{s}(\lambda,x) = (2 \pi i - \log \lambda)^{-s} e^{2\pi i x}$.
Now, observe that, if $s$ is considered fixed, $L$ is $\alpha \odot f_{s}$ as defined in~\eqref{eq:defAction}.


Thus~\eqref{eq:L-inv} is a consequence of~\eqref{eq:invCM}, assuming that the convergence is fast enough. This is so because $\alpha$ and $f_{s}$  satisfy the hypotheses of Theorems~3 and~4 of~\cite{BNV}, which justify the inversion formula. Briefly, since $\alpha(n)$ is a power, all that is needed is the estimate
$d(n) = o(n^{\epsilon})$ for any $\epsilon > 0$, where $d(n)$ is the number of divisors of~$n$.
\end{proof}


Using~\eqref{eq:LW-sep}, we obtain an inversion formula involving the Lerch function. 


\begin{theorem}
\label{th:LW-inv}
For $0 < |\lambda| < 1$, $\Re s > 1$ and $x \in [0,1]$, we have
%
\begin{equation}
\label{eq:LW-inv}
\begin{aligned}
	& (2 \pi i - \log \lambda)^{-s} e^{2 \pi i x}
	  + e^{\pi i s} (2 \pi i + \log \lambda)^{-s} e^{-2 \pi i x}
	\\
	&\quad = \sum_{n=1}^{\infty} \mu(n) \, n^{-s}
	  \left(
	  \frac{\lambda^{\{nx\}/n}}{\Gamma(s)} \Phi(\lambda^{1/n},1-s,\{nx\}) - n^{s} (-\log \lambda)^{-s}
	  \right).
\end{aligned}
\end{equation}
%
\end{theorem}

\begin{proof}
Changing $\lambda$ to $\lambda^{-1}$ and $x$ to $1-x$ in~\eqref{eq:L-inv}, noting that $\{-x\} = 1 - \{x\}$ for $x \notin \Z$, since $L$ is $1$-periodic in $x$, we have
\[
	  (2\pi i + \log\lambda)^{-s} e^{-2\pi i x}
	= \sum_{n=1}^{\infty} \frac{\mu(n)}{n^{s}} L(\lambda^{-1/n},s,1 - \{nx)\}).
\]
Adding this to~\eqref{eq:L-inv} and using~\eqref{eq:LW-sep} we obtain~\eqref{eq:LW-inv}.
\end{proof}

\begin{cor}
\label{cor:AB-inv}
For $k \in \N$, $k \ge 2$, $\lambda \neq 0,1$ and $x \in [0,1]$, we have
%
\begin{equation*}
%\label{eq:AB-inv}
\begin{aligned}
	& (2 \pi i - \log \lambda)^{-k} e^{2 \pi i x}
	  + e^{\pi i k} (2 \pi i + \log \lambda)^{-k} e^{-2 \pi i x}
	\\
	&\qquad = - \sum_{n=1}^{\infty} \mu(n) \, n^{-k}
	  \left(
	  \frac{\lambda^{\{nx\}/n}}{k!} \B_{k}(\{nx\};\lambda^{1/n}) + \frac{n^{k}}{(-\log \lambda)^{k}}
	  \right)
\end{aligned}
\end{equation*}
%
where $\B_{k}(x;\lambda)$ denotes the $k$th Apostol-Bernoulli polynomial.
\end{cor}

\begin{proof}
In the case $0 < |\lambda| < 1$, set $s = k$ for $k \in \N$ in~\eqref{eq:LW-inv} and use~\eqref{eq:A-B}.
Then, since the Apostol-Bernoulli polynomials are defined for $\lambda \neq 0$ and have a unique pole at $\lambda = 1$, the result extends immediately to $\lambda \neq 0,1$ by analytic continuation.
\end{proof}

\begin{remark}
Corollary~\ref{cor:AB-inv} is Theorem~2 of~\cite{Bayad}. There is also a corresponding formula for the Apostol-Euler polynomials $E_{n}(x;\lambda)$ via the simple relation $E_n(x;\lambda) = -\frac{2}{n+1} B_{n+1}(x;-\lambda)$, that was first proved in Lemma~2 of~\cite{NRV-MC}. It should be noted that, although Apostol-Bernoulli and Apostol-Euler polynomials are often dealt with as separate parametrized families, probably due to their specializations to the classical Bernoulli and Euler polynomials, they are essentially one and the same.
\end{remark}

%------------------
\section{Hurwitz's Formula as a limit formula of the Lindel\"of-Wirtinger expansion}
\label{sec:HW}
%------------------


Setting $\lambda = 1$ in the series~\eqref{eq:Lerch} defining the Lerch transcendent $\Phi$ one obtains the definition of the Hurwitz zeta function,
\[
	\zeta(s,z) = \Phi(1,s,z) = \sum_{n=0}^{\infty} \frac{1}{(n + z)^{s}},
\]
where the series converges absolutely	 for $\Re s > 1$ and $z \in \C$, $z \neq 0,-1,-2,\dots$. One often adopts the convention that summands with $n + z = 0$ are omitted. This makes $\zeta(s,0)$ equal by definition to the Riemann zeta function $\zeta(s)$. For $\Re s < 0$, which is the case we consider, one also has $\zeta(s,0) = \zeta(s)$ by analytic continuation.

The Hurwitz zeta function $\zeta(s,z)$ has a well-known analytic continuation to all $s \neq 1$ given by
%
\begin{equation}
\label{eq:HZ-int}
	\zeta(s,x) = \frac{\Gamma(1-s) e^{-\pi i s}}{2\pi i} I(s,x),
	\quad
	I(s,x) = \int_{\mathcal{L}_{\rho}} \frac{z^{s-1} e^{-xz}}{1-e^{-z}} \, dz
\end{equation}
%
where $0 < \rho < 2\pi$ and $\mathcal{L}_{\rho}$ is the path along the positive real axis from $\infty$ to $\rho$, with argument $0$, the counterclockwise circle of radius $\rho$ around the origin, and the path from $\rho$ back to $\infty$, with argument~$2\pi$.

A key relation satisfied by $\zeta(s,z)$ is \emph{Hurwitz's formula}, which is
%
\begin{equation}
	\label{eq:HF}
	\zeta(1-s,x) = \frac{\Gamma(s)}{(2\pi)^{s}}
	               \left( e^{-\pi i s/2} H(s,x) + e^{\pi i s/2} H(s,1-x) \right)
\end{equation}
%
where $\Re s > 1$, $x \in [0,1]$ 
% ($x=0$ is also included if one omits singular terms in the defining series)
and
%
\begin{equation*}
%	\label{eq:periodicZeta}
	H(s,x) = \sum_{n=1}^{\infty} \frac{e^{2 \pi i n x}}{n^{s}}
\end{equation*}
%
is known as the \emph{periodic zeta function.} 

The importance of Hurwitz's formula rests on the fact that at $x = 1$ it specializes to the functional equation of the Riemann zeta function. Many proofs of~\eqref{eq:HF} using different techniques are found in the literature. We will show how Hurwitz's formula can be derived as a limit formula from the Lindel\"of-Wirtinger expansion for the Lerch transcendent and, likewise, M\"obius inversion of Hurwitz's formula from the corresponding inversion of the Lindel\"of-Wirtinger formula given in Theorem~\ref{th:LW-inv}.

Clearly $H(s,x) = (2\pi i)^{s} L(1,s,x)$, where $L$ is defined in~\eqref{eq:L-defi}, and~\eqref{eq:HF} bears an obvious resemblance to~\eqref{eq:LW-sep}.
However, one cannot directly substitute $\lambda = 1$ in formulas such as the Lindel\"of-Wirtinger expansion and obtain a valid result, since $\lambda = 1$ is a logarithmic singularity of $\Phi(\lambda,s,x)$. Rather, the correct approach is to cancel the singularity by subtracting the term corresponding to $n=0$ in~\eqref{eq:L-W}. This is done via an important auxiliary result, namely, the expansion of $\Phi$ in powers of $\log\lambda$. Erd\'elyi showed (see formula~(8) on p.~29 of~\cite{Bate} or formula~(3) in~\cite{Johnson}) that
%
\begin{equation}
\label{eq:Phi-logser}
  \Phi(\lambda,s,x) - \lambda^{-x} \Gamma(1-s){(\log \lambda^{-1})}^{s-1}
  = \lambda^{-x} \sum_{k=0}^{\infty} \zeta(s-k,x) \frac{(\log\lambda)^{k}}{k!}
\end{equation}
%
for $|\log \lambda| < 2 \pi$, $0 < x \leq 1$ and $s \neq 1,2,3,\dots$; for $\Re s < 0$ we may also include $x = 0$.
% This region contains the much smaller disk $|\lambda - 1| \leq \frac{1}{2}$, 
% in which $\log \lambda^{-1} = -\log \lambda$. 

We shall need to justify exchanging the limit as $\lambda \to 1$ with the various infinite sums which appear. This means finding suitable uniform bounds.


\begin{lemma}
\label{lemma:int-bd}
Given $0 < \rho < 2 \pi$ and a compact subset $S$ of the left half-plane $\Re s < 0$, there is a constant $C$ depending only on $S,\rho$, such that the contour integral in~\eqref{eq:HZ-int} satisfies
%
\begin{equation}
\label{eq:int-bd}
	|I(s-k,x)| \leq C \rho^{\sigma-k}
\end{equation}
%
for all $x \in [0,1]$, $s \in S$ and $k \geq 0$
(recall that $\sigma = \Re s$).
\end{lemma}

\begin{proof}
For $\sigma = \Re s < 0$ and $t = \Im s$, we have the estimate $|z^{s-1}| \leq |z|^{\sigma - 1} e^{2 \pi |t|}$, while along the positive real axis we have $|e^{-xz}| = e^{-x \Re z} \leq 1$ and $1 - e^{-z} \geq 1 - e^{-\rho}$. Thus, on either branch along the real axis,
\[
	\left| 
		\int_{\rho}^{\infty} \frac{z^{s-1} e^{-xz}}{1 - e^{-z}} dz
	\right|
	\leq
	  \frac{e^{2\pi |t|}}{1 - e^{-\rho}} \int_{\rho}^{\infty} u^{\sigma - 1} \, du
	= \frac{e^{2\pi |t|}}{1 - e^{-\rho}} \cdot \frac{\rho^{\sigma}}{-\sigma}
\]
and for $|z| = \rho$ we have $|1 - e^{-z}| \geq 1 - e^{\rho - 2\pi}$ and $|e^{-xz}| \leq e^{\rho}$, hence
\[
	\left| 
		\int_{|z|=\rho} \frac{z^{s-1} e^{-xz}}{1 - e^{-z}} dz
	\right|
	\leq	 \frac{2 \pi e^{\rho}  e^{2\pi |t|}}{1 - e^{\rho - 2\pi}} \cdot \rho^{\sigma}.
\]
For any $k \geq 0$, shifting $s$ to $s - k$ with $k \geq 0$, since $k - \sigma \geq -\sigma > 0$, we have $\rho^{\sigma - k}/(k-\sigma) \leq \rho^{\sigma - k}/(-\sigma)$ and hence we obtain the bound
\[
	     |I(s-k,x)|
	\leq   \frac{2e^{2\pi |t|}}{1 - e^{-\rho}} \cdot \frac{\rho^{\sigma-k}}{-\sigma}
	     + \frac{2 \pi e^{\rho}  e^{2\pi |t|}}{1 - e^{\rho - 2\pi}} \cdot \rho^{\sigma-k}.
\]
These bounds clearly imply~\eqref{eq:int-bd} over a compact subset of $\{ \Re s < 0 \}$.
\end{proof}




\begin{prop}
\label{prop:HZ-bd}
Given $0 < \rho < 2\pi$ and a compact subset $S$ of the left half-plane $\Re s < 0$, there is a constant $C$ depending only on $S,\rho$, such that
%
\begin{equation*}
%\label{eq:HZ-bd}
	|\zeta(s-k,x)| \leq C k!\, (2\rho^{-1})^{k}
\end{equation*}
%
for all $x \in [0,1]$ and $k = 0,1,2,\dots$.
\end{prop}


\begin{proof}
We use the integral representation~\eqref{eq:HZ-int}, shifted by~$k$:
\[
	  \zeta(s-k,x) 
%	= \frac{\Gamma(1-s+k) e^{-\pi i(s-k)}}{2\pi i} I(s-k,x)
	= \frac{(-1)^{k} \Gamma(1-s+k) e^{-\pi i s}}{2\pi i} I(s-k,x).
\]
By Lemma~\ref{lemma:int-bd}, we have $|I(s-k,x)| \leq C \rho^{\sigma - k}$ where $C$ depends only on $S,\rho$.
On the other hand, since $\Gamma(1-s+k) = (k-s)(k-s-1)\cdots(-s+1)\Gamma(1-s)$, if $m \in \Z$ is such that $|s| \leq m < |s| + 1$, we have
%%
\begin{align*}
	      \frac{|\Gamma(1-s+k)|}{k!}
%	&\leq \frac{1}{k!}(k + |s|)(k -1 + |s|) \cdots (1 + |s|) |\Gamma(1-s)|
%	\\
	&\leq \frac{1}{k!}(k + m)(k - 1 + m) \cdots (1 + m) |\Gamma(1-s)|
	\\
	&=    \binom{m+k}{k} |\Gamma(1-s)|
	 \leq 2^{m+k} |\Gamma(1-s)|
	 \leq 2^{k + 1 + |s|} |\Gamma(1-s)|.
\end{align*}
Thus, for any $x \in [0,1]$ and $s \in S$, with $\sigma = \Re s$, $t = \Im s$, we have
%%
\begin{align*}
	      \frac{1}{k!}|\zeta(s-k,x)|
	&=    \frac{|\Gamma(1-s+k)| \cdot e^{\pi t}}{2\pi k!} \cdot |I(s-k,x)|
	\\
	&\leq \pi^{-1} 2^{k+|s|} \cdot |\Gamma(1-s)| \cdot e^{\pi |t|} \cdot C\rho^{\sigma-k}
	 \leq C' (2 \rho^{-1})^{k}
\end{align*}
%%
where $C'$ is another constant depending only on $S,\rho$.
\end{proof}


\begin{prop}
\label{prop:logser-unifconv}
The tails of the series~\eqref{eq:Phi-logser} satisfy a bound of the form
%
\begin{equation}
\label{eq:tail-bd}
     \sum_{k=m}^{\infty} \left| \zeta(s-k,x) \frac{\log^{k}\lambda}{k!}\right|
\leq C |\log\lambda|^{m}
\end{equation}
%
for $x \in [0,1]$ and $\lambda,s$ in respective compact subsets $\Lambda,S$ of the domains $|\log\lambda| < \pi$ and $\Re s < 0$, where $C > 0$ is a constant depending only on $\Lambda,S$. Thus~\eqref{eq:Phi-logser} is normally convergent for $x \in [0,1]$ and on compact subsets of $\Re s < 0$ and $|\log\lambda| < \pi$.
\end{prop}

\begin{proof}
By Proposition~\ref{prop:HZ-bd}, given $0 < \rho < 2\pi$ and a compact subset $S$ of the left half-plane $\Re s < 0$, there is a constant $C(S,\rho)$ depending only on $\rho$ and~$S$, such that $\frac{1}{k!}|\zeta(s-k,x)| \leq C(S,\rho) (2\rho^{-1})^{k}$ for all $s \in S$, $x \in [0,1]$, and $k \geq 0$.

For $\lambda$ in a compact subset $\Lambda$ of the domain $|\log\lambda|<\pi$, there is a uniform bound $|\log \lambda| \leq \ell < \pi$. Since $0 < \rho < 2\pi$ is arbitrary, choosing $2 \ell < \rho < 2\pi$ with $\rho$ close enough to $2\pi$ makes $2 \rho^{-1} \ell < 1$ and $2\rho^{-1} < 1$, hence for $(\lambda,s,x) \in \Lambda \times S \times [0,1]$, the series~\eqref{eq:tail-bd} is dominated by a convergent geometric series, yielding the bound
\[
	 \sum_{k=m}^{\infty} \frac{1}{k!}|\zeta(s-k,x)| \, |\log\lambda|^{k}
	 \leq \frac{C(S,\rho)}{1 - 2\rho^{-1}\ell} \, {|2\rho^{-1}\log\lambda|}^{m}
	 \leq C(S,\Lambda) |\log\lambda|^{m}
\]
for all $m \geq 0$, where $C(S,\Lambda) > 0$ is a constant depending only on $S,\Lambda$.
\end{proof}

\begin{cor}
\label{cor:Hardy-rel}
For $x \in [0,1]$ and $\lambda,s$ lying in respective compact subsets $\Lambda,S$ of the domains $|\log \lambda| < \pi$, $0 < |\lambda| < 1$, and $\Re s < 0$, there is a constant $C > 0$ depending only on $\Lambda,S$ such that
%
\begin{equation}
\label{eq:Hardy-bd}
	\big|\Phi(\lambda,s,x) - \lambda^{-x}\Gamma(1-s){(\log \lambda^{-1})}^{s-1} - \lambda^{-x}\zeta(s,x)\big| 
	\leq C |\log \lambda|
\end{equation}
%
for $(\lambda,s,x) \in \Lambda \times S \times [0,1]$.
In particular, \emph{Hardy's relation} (formula (4) in~\cite{Johnson}) holds uniformly:
%
\begin{equation}
\label{eq:Hardy-rel}
	\lim_{\lambda \to 1} \left( \Phi(\lambda,s,x) - \lambda^{-x} \Gamma(1-s){(\log \lambda^{-1})}^{s-1}\right) 
	= \zeta(s,x)
\end{equation}
%
for $(s,x) \in S \times [0,1]$.
\end{cor}

\begin{proof}
Assume $s$ lies in a compact subset $S$ of $\{ \Re s < 0 \}$ and $\lambda$ in a compact subset $\Lambda$ of $|\log \lambda| < \pi$. Separating the summand corresponding to $k = 0$ in Erd\'elyi's expansion~\eqref{eq:Phi-logser}, the bound~\eqref{eq:tail-bd} of Proposition~\ref{prop:logser-unifconv} shows that
\[
	\big|\Phi(\lambda,s,x) - \lambda^{-x}\Gamma(1-s){(\log \lambda^{-1})}^{s-1} - \lambda^{-x}\zeta(s,x)\big| 
%	= \lambda^{-x} \sum_{k=1}^{\infty} \zeta(s-k,x) \frac{\log^{k}\lambda}{k!}.
	\leq C |\lambda|^{-x} |\log \lambda|
\]
where $C > 0$ is a constant depending only on $S,\Lambda$. Now, $\min(1,|\lambda|) \leq |\lambda^{x}| \leq \max(1,|\lambda|)$ for $x \in [0,1]$, and the mean value inequality for the exponential gives $|e^{z}-1| \leq |z| e^{\Re^{+} z}$ where $\Re^{+} = \max(0,\Re)$. Thus we have
\[
    |\lambda^{x} - 1| \leq |x \log\lambda| e^{x \Re^{+} \log\lambda} \leq |\log\lambda|\cdot \max(1,|\lambda|)
\]
for $x \in [0,1]$ and hence, since $\zeta(s,x)$ is bounded on $S \times [0,1]$,
\begin{multline*}
\qquad
	\big|\Phi(\lambda,s,x) - \lambda^{-x}\Gamma(1-s){(\log \lambda^{-1})}^{s-1} - \zeta(s,x)\big| 
	\\
  \begin{aligned}
	&\leq C |\lambda|^{-x} |\log \lambda| + |\lambda^{-x} - 1| |\zeta(s,x)|
	\\
	&\leq C |\lambda|^{-x} |\log \lambda| + |\lambda|^{-x}|\lambda^{x} - 1| \cdot \|\zeta\|_{S \times [0,1]}
%	\\
%	&\leq C |\lambda|^{-x} |\log \lambda| + |\lambda|^{-x} |\log\lambda| \max(1,|\lambda|) \cdot \|\zeta\|_{S \times [0,1]}
	\\
	&\leq \frac{C  + \max(1,|\lambda|) \cdot \|\zeta\|_{S \times [0,1]}}{\min(1,|\lambda|)} \,|\log\lambda| 
%	\\
	\leq C' |\log\lambda|
  \end{aligned}
\qquad
\end{multline*}
for another constant $C'$ depending only on $S,\Lambda$. Finally, note that $|\log \lambda| = O(|\lambda-1|)$ as $\lambda \to 1$.
\end{proof}

\begin{cor}
\label{cor:LW->HF}
The Lindel\"of-Wirtinger expansion~\eqref{eq:L-W} implies Hurwitz's formula~\eqref{eq:HF}.
\end{cor}

\begin{proof}
By Hardy's relation~\eqref{eq:Hardy-rel}, changing $s$ to $1-s$ and hence assuming $\Re s > 1$, we have
\[
	\lim_{\lambda \to 1} \left(\Phi(\lambda,s,x) - \lambda^{-x}\Gamma(s){(\log \lambda^{-1})}^{-s}\right) = \zeta(1-s,x).
\]
Now take the limit as $\lambda \to 1$ in the modified form~\eqref{eq:LW-sep} of the Lindel\"of-Wirtinger relation, noting that $\log \lambda^{-1} = -\log \lambda$ for $0 < |\lambda| < 1$:
%
\begin{multline*}
  \qquad\qquad
	  \Phi(\lambda,1-s,x) - \lambda^{-x}\Gamma(s){(\log \lambda^{-1})}^{-s} \\
	  = \lambda^{-x}\Gamma(s)\left(L(\lambda,s,x) + e^{\pi i s} L(\lambda^{-1},s,1-x)\right)
  \qquad\qquad
\end{multline*}
%
to obtain, since $L(\lambda,s,x)$ is holomorphic for $\lambda \notin (-\infty,0]$,
%%
\begin{align*}
	   \zeta(1-s,x) 
	&= \Gamma(s)\left(L(1,s,x) + e^{\pi i s}L(1,s,1-x)\right)
	\\
	&= \Gamma(s)(2\pi i)^{-s}\left(H(s,x) + e^{\pi i s}H(s,1-x)\right)
%	\\
%	&= \frac{\Gamma(s)}{(2\pi)^{s}} e^{-\pi i s/2}\left(H(s,x) + e^{\pi i s}H(s,1-x)\right)
	\\
	&= \frac{\Gamma(s)}{(2\pi)^{s}} \left(e^{-\pi i s/2}H(s,x) + e^{\pi i s/2}H(s,1-x)\right),
\end{align*}
%%
which is Hurwitz's formula.
\end{proof}


%------------------
\section{M\"obius inversion of Hurwitz's formula}
\label{sec:HW-Inv}
%------------------


One can apply the method of M\"obius inversion outlined in Section~\ref{sec:MobiusInvGen} to Hurwitz's formula~\eqref{eq:HF} in the same manner as it was used to invert the Lindel\"of-Wirtinger expansion in Section~\ref{sec:LW-inv}. Rather than repeating the process, we may derive the result as a limit formula of~\eqref{eq:LW-inv} when $\lambda \to 1$, by subtracting the logarithmic singularity using the technique outlined in the previous section.


\begin{prop}
\label{prop:LW-inv-conv}
The series~\eqref{eq:LW-inv} resulting from M\"obius inversion of the Lindel\"of-Wirtinger expansion in the form~\eqref{eq:LW-sep}, is normally convergent for $x \in [0,1]$, $s$ in a compact subset of the half-plane $\Re s > 1$, and $\lambda$ in a compact subset of the domain $|\log\lambda|<\pi$, $0 < |\lambda| < 1$. The following limit formula holds:
%
\begin{equation}
\label{HF-inv}
	e^{-\pi i s/2} e^{2 \pi i x} + e^{\pi i s/2} e^{-2 \pi i x}
	= \frac{(2\pi)^{s}}{\Gamma(s)} \sum_{n=1}^{\infty} \frac{\mu(n)}{n^{s}} \zeta(1-s,\{nx\})
\end{equation}
%
where $\{x\}$ denotes the fractional part of~$x$.
\end{prop}


\begin{proof}
Changing $s$ to $1-s$ in the bound~\eqref{eq:Hardy-bd} which implies Hardy's relation~\eqref{eq:Hardy-rel}, we have
\[
	\big|\Phi(\lambda,1-s,x) - \lambda^{-x}\Gamma(s){(\log \lambda^{-1})}^{-s} - \zeta(1-s,x)\big| \leq C |\log\lambda|
\]
for $x \in [0,1]$ and $\lambda,s$ in respective compact subsets $\Lambda,S$ of $|\log\lambda|<\pi$ and $\Re s > 1$, where $C > 0$ depends only on $\Lambda,S$.
Given $n = 1,2,3,\dots$, changing $\lambda$ to $\lambda^{1/n}$ (the principal branch) and $x$ to the fractional part $\{nx\}$, we obtain, assuming also that $\lambda \notin (-\infty,0]$,
%
\begin{multline*}
	\big|\Phi(\lambda^{1/n},1-s,\{nx\}) - \lambda^{-\{nx\}/n}\Gamma(s) \, n^{s} {(\log \lambda^{-1})}^{-s} - \zeta(1-s,\{nx\})\big|
	\\ 
	\leq C n^{-1} |\log\lambda|
	\leq C |\log\lambda|,
\end{multline*}
%
so that the limit as $\lambda \to 1$ is uniform in $s,x$.
Now, rewrite~\eqref{eq:LW-inv} as
%
\begin{multline*}
	(2 \pi i - \log \lambda)^{-s} e^{2 \pi i x}
	  + e^{\pi i s} (2 \pi i + \log \lambda)^{-s} e^{-2 \pi i x}
	\\
	= \frac{1}{\Gamma(s)}
	  \sum_{n=1}^{\infty} \frac{\mu(n)}{n^{s}} \, \lambda^{\{nx\}/n}
	  \left(
		\Phi(\lambda^{1/n},1-s,\{nx\}) - \lambda^{-\{nx\}/n} \Gamma(s) n^{s} (-\log \lambda)^{-s}
	  \right)
\end{multline*}
%
and note that $|\lambda^{\{nx\}/n}| \leq \max(1,|\lambda|)$, so that the series 
\[
	\frac{1}{\Gamma(s)} \sum_{n=1}^{\infty} \frac{\mu(n)}{n^{s}} \, \lambda^{\{nx\}/n},
	\quad
	\frac{1}{\Gamma(s)} \sum_{n=1}^{\infty} \frac{\mu(n)}{n^{s}} \, \lambda^{\{nx\}/n} \zeta(1-s,\{nx\})
\]
are both normally convergent for $(\lambda,s,x) \in \Lambda \times S \times [0,1]$, and hence taking the limit as $\lambda \to 1$ in~\eqref{eq:LW-inv} yields
\[
	(2\pi i)^{-s} \left(e^{2\pi i x} + e^{\pi i s} e^{-2\pi i x}\right)
	= \frac{1}{\Gamma(s)} \sum_{n=1}^{\infty} \frac{\mu(n)}{n^{s}} \, \zeta(1-s,\{nx\})
\]
which is equivalent to~\eqref{HF-inv}.
\end{proof}


\begin{remark}
Note that~\eqref{HF-inv} has the equivalent form
%
\begin{equation}
\label{eq:HF-inv-cos}
	  2 (2 \pi)^{-s} \Gamma(s) \cos\left(\pi\left(\frac{s}{2} - 2 x\right)\right)
	= \sum_{n=1}^{\infty} \frac{\mu(n)}{n^{s}} \, \zeta(1-s,\{nx\}).
\end{equation}
%
\end{remark}


\begin{cor}
\label{cor:Ber-inv}
For $k = 2,3,4,\dots$ and $x \in [0,1]$, we have
%
\begin{equation*}
%\label{eq:Ber-inv}
	  \cos\left( \frac{\pi k}{2} - 2 \pi x \right)
	= - \frac{(2\pi)^{k}}{2 (k!)} \sum_{n=1}^{\infty} \frac{\mu(n)}{n^{k}} B_{k}(\{nx\}),
\end{equation*}
%
where $B_{k}(x)$ is the $k$th Bernoulli polynomial. Equivalently, separating according to the parity of $k$, we obtain the pair of formulas ((2.4) and (2.5) of~\cite{NRV-JAT})
%
\begin{equation*}
%	\label{eq:Ber-inv-sep}
	\begin{aligned}
		\cos(2 \pi x) &= (-1)^{k-1} \frac{(2\pi)^{2k}}{2 (2k)!}
		                 \sum_{n=1}^{\infty} \frac{\mu(n)}{n^{2k}} B_{2k}(\{nx\}),
					  \\
		\sin(2 \pi x) &= (-1)^{k-1} \frac{(2\pi)^{2k+1}}{2 (2k+1)!}
		                 \sum_{n=1}^{\infty} \frac{\mu(n)}{n^{2k+1}} B_{2k+1}(\{nx\}).
	\end{aligned}
\end{equation*}
%
\end{cor}

\begin{proof}
Let $s = k$ with $k \in \N$, $k > 1$, in~\eqref{HF-inv}.
The Hurwitz zeta function then evaluates to
%
\begin{equation*}
%	\label{eq:HZ-ber}
	\zeta(1-k,x) = - \frac{B_{k}(x)}{k}
\end{equation*}
%
for $0 \leq x \leq 1.$
\end{proof}



%------------------
\section{A Functional Equation}
\label{sec:functionaleq}
%------------------


Since $\sum_{n=1}^{\infty} \frac{\mu(n)}{n^{s}} = \zeta(s)^{-1}$, for $x = 1$ and $\Re s > 1$, the expansion~\eqref{HF-inv} and its equivalent~\eqref{eq:HF-inv-cos} reduce to the functional equation of the Riemann zeta function, in the form
\[
	\zeta(1-s) = 2 (2 \pi)^{-s} \Gamma(s) \cos\left(\frac{\pi s}{2}\right) \zeta(s)
\]
(using the reflection formula for the gamma function yields a similar relation but with the sine function).


In general, \eqref{HF-inv} may be regarded as a combination of the functional equation with the discrete dynamical system of fractional parts of multiples of a real number, which is conjugate to the rotations of the circle. Recall that $\{nx\}$ is equidistributed in $[0,1]$ when $x$ is irrational, while it is a periodic sequence when $x$ is rational. 

For example, consider the Dirichlet series
%
\begin{equation*}
%	\label{eq:Zinv-mod}
		  \zeta^{*}(m,r,s)
		= \sum_{n \equiv r \bmod m} \frac{\mu(n)}{n^{s}},
		\qquad
		\Re s > 1,
\end{equation*}
%
for $r,m \in \N$ with $1 \leq r \leq m$.
If $x \in \Q$ has denominator equal to $m$, then $\{nx\}$ is $m$-periodic and for $\Re s > 1$ we can group terms in~\eqref{eq:HF-inv-cos} according to their residue modulo~$m$. We obtain
%
\begin{equation*}
%	\label{eq:HF-inv-cos-rat}
%	\begin{aligned}
	   2 (2 \pi)^{-s} \Gamma(s) \cos\left(\pi\left(\frac{s}{2} - 2 x\right)\right)
%	&=  \sum_{n=1}^{\infty} \frac{\mu(n)}{n^{s}} \zeta(1-s,\{nx\})
%	\\
	= \sum_{r=1}^{m} \zeta^{*}(m,r,s) \zeta(1-s,\{rx\}),
%	\end{aligned}
\end{equation*}
%
which again reduces to the functional equation for $\zeta(s)$ for $r = m = 1$.


%:bibliography
%--------------------------------------
\begin{thebibliography}{00}
%--------------------------------------

\bibitem{Ap}
T. M. Apostol,
On the Lerch zeta function,
\textit{Pacific J. Math.}
\textbf{1} (1951), 161--167;
addendum, \textit{Pacific J. Math.} \textbf{2} (1952), 10.

\bibitem{Ap-NIST}
T. M. Apostol,
Zeta and related functions,
\textit{NIST handbook of mathematical functions}
(edited by F.~W. F.~Olver, D.~W. Lozier, R.~F. Boisvert and C.~W. Clark),
601--616,
National Institute of Standards and Technology, Washington, DC, and
Cambridge University Press, Cambridge, 2010.
Available online in \url{http://dlmf.nist.gov/25}

\bibitem{BaSa}
E. P. Balanzario and J. S\'{a}nchez-Ortiz,
Riemann-Siegel integral formula for the Lerch zeta function,
\textit{Math. Comp.}
\textbf{81} (2012), 2319--2333.
% DOI: \url{http://dx.doi.org/10.1090/S0025-5718-2011-02566-4}

\bibitem{Bayad}
A. Bayad and J. Chikhi,
M\"obius inversion formulae for Apostol-Bernoulli type polynomials and numbers,
\textit{Math. Comp.}
\textbf{82} (2013), 2327--2332.

\bibitem{BNV}
M. Benito, L. M. Navas and J. L. Varona,
M\"{o}bius inversion formulas for flows on arithmetic semigroups,
\textit{Journal of Number Theory}
\textbf{128} (2008), 390--412.


\bibitem{BNV2}
M. Benito, L. M. Navas and J. L. Varona,
M\"{o}bius inversion from the point of view of arithemtical semigroup flows, in
Proceedings of the ``Segundas Jornadas de Teor\'{\i}a de N\'{u}meros'', Madrid (2007),
\textit{Bibl. Rev. Mat. Iberoamericana}
(2008), 63--81.


\bibitem{Boy}
K. N. Boyadzhiev,
Apostol-Bernoulli functions, derivative polynomials and Eulerian polynomials,
\textit{Adv. Appl. Discrete Math.}
\textbf{1} (2008), 109--122.

\bibitem{Bate} % Bateman project
A. Erd\'{e}lyi, W. Magnus, F. Oberhettinger and F. Tricomi,
\textit{Higher Transcendental Functions}, Volume I,
McGraw-Hill, New York, 1953.

\bibitem{GuiS}
J. Guillera and J. Sondow,
Double integrals and infinite products for some classical constants via analytic continuations of Lerch's transcendent,
\textit{Ramanujan J.}
\textbf{16} (2008), 247--270.
% Preprint in \url{http://arxiv.org/abs/math.NT/0506319}

\bibitem{Johnson}
B. R. Johnson,
Generalized Lerch Zeta Function,
\textit{Pac. J. Math.}
\textbf{53} (1974), 189--193.

\bibitem{Lerch}
M. Lerch,
Note sur la fonction $\mathfrak{K}(w,x,s) = \sum_{k=0}^{\infty} e^{2\pi ikx}(w+k)^{-s}$,
\textit{Acta Math.}
\textbf{11} (1887), 19--24.
% \url{http://www.springerlink.com/content/bpn8104100w31667/}
% http://www.springerlink.com/content/0001-5962/11/1-4/

\bibitem{NRV-JAT}
L. M. Navas, F. J. Ruiz and J. L. Varona,
The M\"obius inversion formula for Fourier series
applied to Bernoulli and Euler polynomials,
\textit{J. Approx. Theory}
\textbf{163} (2011), 22--40.

\bibitem{NRV-MC}
L. M. Navas, F. J. Ruiz and J. L. Varona,
Asymptotic estimates for Apostol-Bernoulli and Apostol-Euler polynomials,
\textit{Math. Comp.}
\textbf{81} (2012), 1707--1722.

\bibitem{NRV-JAT2}
L. M. Navas, F. J. Ruiz and J. L. Varona,
Asymptotic behavior of the Lerch transcendent function,
\textit{J. Approx. Theory}
\textbf{170} (2013), 21--31. 
% DOI: \url{http://dx.doi.org/10.1016/j.jat.2012.03.006}

\bibitem{SC-book}
H. M. Srivastava and J. Choi,
\textit{Series associated with the zeta and related functions},
Kluwer Academic, Dordrecht, 2001.

\bibitem{Wir}
W. Wirtinger,
\"Uber eine besondere Dirichletsche Reihe,
\textit{J. Reine Angew. Math.} \textbf{129} (1905), 214--219.
% % Journal f\"ur die Reine und Angewandte Mathematik

%--------------------------------------
\end{thebibliography}
%--------------------------------------

%--------------------------------------
\end{document}
%--------------------------------------
